\setcounter{ex}{0}
\setcounter{paragraph}{0}
\PhanI
%\loai{Xác định biên độ dao động và tần số góc}

\begin{ex} %[3P1Y1-2][Loai1] 
(QG - 2016) Một chất điểm dao động có phương trình $x=10\cos\left(15t+\pi\right)$ (x tính bằng\ cm, t tính bằng s). Chất điểm này dao động với tần số góc là
\choice
{$20\ rad/s$}
{$10\ rad/s$}
{$5\ rad/s$}
{\True $15\ rad/s$}
\loigiai{
Ta có: $\omega =15\ rad/s$.
}
%\haidong{2}
\end{ex} \haidong{15}
\begin{ex}
Một chất điểm dao động điều hoà có phương trình li độ theo thời gian là $x=6\cos(4\pi t+\dfrac{\pi}{3})\ (cm)$. Chu kì của dao động là
\choice
{4 s}
{2 s}
{0,25 s}
{\True 0,5 s}
\loigiai{
\begin{itemize}
\item Chu kì của dao động bằng: $T=\dfrac{2\pi}{\omega}=0,5\ s$
\end{itemize}
}
%\haidong{2}
\end{ex} \haidong{15}
\begin{ex}
Một chất điểm dao động điều hoà có tần số góc $\omega=10\pi\ rad/s$. Tần số của dao động là
\choice
{\True $5\ Hz$}
{$10\ Hz$}
{$20\ Hz$}
{$5\pi\ Hz$}
\loigiai{
\begin{itemize}
\item Tần số của dao động là: $f=\dfrac{\omega}{2\pi}=5\ Hz$.
\end{itemize}
}
%\haidong{3}
\end{ex} \haidong{15}
%\loai{Chiều dài quỹ đạo}
\begin{ex}%[3P1B1-2][Loai1]
Một chất điểm dao động theo phương trình $x = 6\cos\omega t$ cm. Quỹ đạo chất điểm có độ dài bằng
\choice
{2 cm}
{6 cm}
{3 cm}
{\True 12 cm}
\loigiai{Quỹ đạo chất điểm có độ dài: $L=2A=12\ cm$.
}
%\haidong{3}
\end{ex} \haidong{15}

\begin{ex}%[3P1Y1-2]
Một vật nhỏ dao động điều hòa theo một quỹ đạo dài 24 cm. Dao động có biên độ
\choice
{6 cm}
{3 cm}
{\True 12 cm}
{24 cm}
\loigiai{
Quỹ đạo của dao động điều hòa bằng: $L = 2A =24\ cm \Rightarrow A = 12\  cm$.
}
%\haidong{2}
\end{ex} \haidong{15}

 
\begin{ex} %[3P1B1-2][Loai1]  
Một chất điểm dao động tuần hoàn với chu kì 0,2 s. Biết rằng trong 2 s, chất điểm đi được tổng quãng đường dài 1,6 m. Độ dài quỹ đạo chuyển động của chất điểm là
\choice
{4 cm}
{\True 8 cm}
{12 cm}
{16 cm}
\loigiai{
\begin{itemize}
\item Ta có: $2\ s = 10.0,2 = 10T \Rightarrow S= 10.4A = 160 \Rightarrow A = 4\ cm$.
\item Độ dài quỹ đạo là $L = 2A = 8$ cm.
\end{itemize}}
%\haidong{2}
\end{ex} \haidong{15}
%\loai{Quãng đường đi trong một chu kì}

\begin{ex}%[3P1-1.2B][Loai1]  
Một vật dao động điều hòa đi được quãng đường 16 cm trong một chu kì dao động. Biên độ dao động của vật là
\choice
{\True 4 cm}
{8 cm}
{16 cm}
{2 cm}
\loigiai{
Quãng đường mà vật đi được trong một chu kì:
$S=4A=16\ cm\Rightarrow A=4\ cm.$
}
%\haidong{2}
\end{ex} \haidong{15}
%\loai{Pha của dao động}
\begin{ex}%[3P1Y1-2]
Một chất điểm dao động theo phương trình $x = 10\cos2\pi t$ cm có pha tại thời điểm t là
\choice
{\True $2\pi t$}
{0}
{$2\pi$}
{$\pi$}
\loigiai{
Phương trình dao động tổng hợp là: $x = A\cos \left(\omega t + \varphi\right)$ cm $\Rightarrow$ pha tại thời điểm t là $2\pi t$.
}
%\haidong{2}
\end{ex} \haidong{15}
%\loai{Xác định li độ}
\begin{ex}
Một chất điểm dao động điều hoà với phương trình $x=5\cos(10\pi t+\dfrac{\pi}{3})\ (cm)$. Li độ của vật khi pha dao động bằng $\pi$ là
\choice
{$5\ cm$}
{\True $-5\ cm$}
{$2,5\ cm$}
{$-2,5\ cm$}
\loigiai{
Li độ của vật khi pha dao động bằng $\pi$ là $-5\ cm$.
}
%\haidong{2}
\end{ex} \haidong{15}
\begin{ex}
Một chất điểm dao động điều hoà có phương trình li độ theo thời gian là  $x=5\sqrt 3\cos(10\pi t+\dfrac{\pi}{3})\ (cm)$. Tại thời điểm $t=1\ s$ thì li độ của vật bằng
\choice
{$2,5\ cm$}
{$-5\sqrt 3\ cm$}
{$5\ cm$}
{\True $2,5\sqrt 3\ cm$}
\loigiai{
Tại thời điểm $t=1\ s$ thì li độ của vật bằng $2,5\sqrt 3\ cm$.
}
%\haidong{2}
\end{ex} \haidong{15}
%\loai{Tính chu kì, tần số}
\begin{ex}%[3P1B1-2]
Một vật nhỏ dao động điều hòa thực hiện 2016 dao động toàn phần trong 1008 s. Tần số dao động là
\choice
{$4\pi$  Hz}
{0,5 Hz}
{\True 2 Hz}
{1 Hz}
\loigiai{
Tần số dao động là $f=\dfrac{2016}{1008}=2 \ Hz.$
}
%<MyLT1>
%\haidong{2}
\end{ex} \haidong{15}

\begin{ex}
Một chất điểm dao động điều hoà. Trong thời gian 1 phút vật thực hiện được 30 dao động. Chu kì dao động của vật là
\choice
{\True 2 s}
{30 s}
{0,5 s}
{1 s}
\loigiai{
\begin{itemize}
\item $T=\dfrac{\Delta t}{N}=2\ s$.
\end{itemize}
}
%\haidong{3}
\end{ex} \haidong{15}


%\loai{Xác định pha ban đầu của hàm cos}

%\loai{Xác định pha ban đầu của hàm -cos}
\begin{ex}
Phương trình dao động của một vật có dạng: $x=-A\cos(\omega t+\dfrac{\pi}{3})\ (cm)$. Pha ban đầu của dao động là
\choice
{$\dfrac{\pi}{3}$}
{$-\dfrac{\pi}{3}$}
{$\dfrac{2\pi}{3}$}
{\True $-\dfrac{2\pi}{3}$}
\loigiai{
\begin{itemize}
\item $\varphi=\dfrac{\pi}{3}-\pi=-\dfrac{2\pi}{3}$.
\end{itemize}
}%\haidong{2}
\end{ex} \haidong{15}
%\loai{Xác định pha ban đầu của hàm sin}
\begin{ex} %[3P1B1-2][Loai1]  
Phương trình dao động điều hòa của một chất điểm có dạng: $x = 5\sin\left(6t + \dfrac{\pi}{3}\right)$ cm. Pha ban đầu của dao động trên là
\choice
{\True $-\dfrac{\pi}{6}$ rad}
{$\dfrac{\pi}{6}$ rad}
{$\dfrac{\pi}{3}$ rad}
{$-\dfrac{\pi}{3}$ rad}
\loigiai{
\begin{itemize}
\item Ta có thể viết: $x = 5\sin\left(6t + \dfrac{\pi}{3}\right) = 5\cos\left(6t + \dfrac{\pi}{3} - \dfrac{\pi}{2}\right) = 5\cos\left(6t -\dfrac{\pi}{6}\right)$ cm.
\item Pha ban đầu của dao động là $-\dfrac{\pi}{6}$ rad.
\end{itemize}
}
%\haidong{2}
\end{ex} \haidong{15}
\begin{ex} %[3P1B1-2][Loai1]  
Một dao động điều hòa với phương trình $x = 4\sin\left(\pi t + \dfrac{\pi}{3}\right)$ cm. Pha dao động ban đầu  của vật là
\choice
{$\dfrac{\pi}{6}$ rad}
{$-\dfrac{\pi}{3}$ rad}
{$\dfrac{\pi}{3}$ rad}
{\True $-\dfrac{\pi}{6}$ rad}
\loigiai{
\begin{itemize}
\item Ta có:
$x = 4\sin\left(\pi t+ \dfrac{\pi}{3}\right)\ cm= 4\cos\left(\pi t+ \dfrac{\pi}{3}- \dfrac{\pi}{2}\right)\ cm = 4\cos\left(\pi t- \dfrac{\pi}{6}\right)$ cm.
\item Pha ban đầu của dao động trên là $-\dfrac{\pi}{6}$ rad.
\end{itemize}}
%<MyLT1>
%\haidong{2}
\end{ex} \haidong{15}

%\loai{Xác định pha ban đầu khi \bth{x=A\cos(\omega t+\varphi)} với $\varphi>\pi$ hoặc $\varphi<-\pi$}
\begin{ex} %[3P1B1-2][Loai1]  
Một vật dao động điều hòa với phương trình: $x = 8\cos\left(4\pi t + \dfrac{4\pi}{3}\right)$ cm. Pha ban đầu của dao động của vật là
\choice
{$-\dfrac{\pi}{3}$ rad}
{$\dfrac{\pi}{3}$ rad}
{\True $-\dfrac{2\pi}{3}$ rad}
{$\dfrac{2\pi}{3}$ rad}
\loigiai{
\begin{itemize}
\item Ta có thể viết: $x = 8\cos\left(8\pi t + \dfrac{4\pi}{3}\right) = 8\cos\left(8\pi t - \dfrac{2\pi}{3}\right)\ cm$. Pha ban đầu của dao động là $ \dfrac{2\pi}{3} $ rad.
\end{itemize}
}
%\haidong{2}
\end{ex} \haidong{15}
\begin{ex} %[3P1B4-1][Loai1]  
Một dao động điều hòa với phương trình $x = 4\cos\left(2\pi t + \dfrac{6\pi}{5}\right)$ cm. Pha ban đầu của dao động của vật là
\choice
{$-\dfrac{\pi}{5}$ rad}
{$\dfrac{4\pi}{5}$ rad}
{\True $-\dfrac{4\pi}{5}$ rad}
{$\dfrac{\pi}{5}$ rad}
\loigiai{
\begin{itemize}
\item Ta có $x = 4\cos\left(2\pi t+ \dfrac{6\pi}{5}\right)\ cm= 4\cos\left(2\pi t + \dfrac{6\pi}{5}-2\pi\right)\ cm= 4\cos\left(2\pi t- \dfrac{4\pi}{5}\right)$ cm. Pha ban đầu của dao động của vật là $-\dfrac{4\pi}{5}$ rad.
\end{itemize}}

%\haidong{2}
\end{ex} \haidong{15}
\begin{ex}%[3P1-19.1K][Loai1]
\immini[thm]{
Một chất điểm dao động điều hòa có đồ thị biểu diễn sự phụ thuộc của li độ x theo thời gian t như hình bên. Tần số dao động của chất điểm bằng
\choice
{$0,5\pi$ rad/s}
{$0,5$ Hz}
{$\pi$ rad/s}
{\True $0,25$ Hz}}{
\pgfplotsset{width=7cm,height=4cm,compat=1.4}
\begin{tikzpicture}[draw=Mapcolor,very thick,>=stealth']
	\pgfplotsset{width=7cm,height=4cm,compat=1.4}
	\begin{axis}[
	axis lines=middle, xmin=0, xmax=2.2*pi, xstep=1,
	ymin=-4.5, ymax=5, ystep=0.1,
	grid=major,%Vẽ chế độ lưới theo trục tọa độ Oxy,
	x grid style={dashed, Mapcolor},
	xlabel=$t~(s)$,
	xlabel style={above},
	axis line style={very thick,Mapcolor},
	xtick ={1/1,2/1,3/1,4/1,5/1,6/1,7/1,8/1,9/1},
	xticklabels={1,2},
	y grid style={dashed, Mapcolor},
	ylabel=$x$,
	ylabel style={above},
	ytick={-4, -2, 0, 2, 4},
	yticklabels={},]
	\addplot[line width=1.2pt, domain=0:2*pi, samples=400, Mapcolor]{4*cos(deg(0.5*pi*x+0))};
	\addplot[Mapcolor,mark=*,mark options={fill=white,mark size=1.2pt}, nodes near coords={}]
	coordinates {(1, {4*cos(deg(0.5*pi*x+0))})};
	\addplot[Mapcolor,mark=*,mark options={fill=white,mark size=1.2pt}, nodes near coords={}]
	coordinates {(2, 0)};
	\end{axis}
\end{tikzpicture}
}
\loigiai{
Từ đồ thị, ta thu được: $T=4\ s\Rightarrow f=\dfrac{1}{T}=0,25\ Hz$.
}
%\haidong{3}
\end{ex} \haidong{20}
\setcounter{ex}{0}
\PhanII
\begin{ex}
Một vật đang dao động điều hòa theo phương trình $x=-5 \cos 5 \pi t\ cm$.
\choiceTF[t]
{Pha ban đầu của li độ bằng 0}
{\True Biên độ dao động là $5 \ cm$}
{\True Tần số dao động là $2,5 \ Hz$}
{\True Chu kì dao động là $0,4 \ s$}
\loigiai{
\begin{itemchoice}
\itemch
\itemch
\itemch
\itemch
\end{itemchoice}
}
%\haidong{3}
\end{ex} \haidong{15}
\begin{ex}
Một chất điểm dao động điều hòa với phương trình $x=2 \cos (4 \pi t-0,5 \pi)\ cm$.
\choiceTF[t]
{Pha dao động của chất điểm là $0,5 \pi\ rad$}
{Tần số góc của chất điểm là $2\ rad/s$}
{Chu kì dao động của chất điểm là $1 \ s$}
{\True Chất điểm dao động trên quỹ đạo có chiều dài $4 \ cm$}
\loigiai{
\begin{itemchoice}
\itemch
\itemch
\itemch
\itemch
\end{itemchoice}
}
\end{ex} \haidong{15}
\begin{ex}
Một vật dao động điều hòa trên trục Ox với phương trình $x=4 \cos \left(2 \pi t-\dfrac{\pi}{6}\right)\ (cm)$, (trong đó t tính bằng giây).
\choiceTF[t]
{\True Biên độ dao động của vật là $4\ cm$}
{Pha ban đầu của dao động là $\dfrac{\pi}{6}$ rad}
{Chiều dài quỹ đạo chuyển động của vật là $4\ cm$}
{\True Tại thời điểm $t=1 \ s$, vật cách vị trí cân bằng một đoạn $2 \sqrt{3} \ cm$}
\loigiai{
\begin{itemchoice}
\itemch Biên độ là hệ số của cos, $A = 4$ cm.
\itemch Pha ban đầu là $-\pi/6$ rad.
\itemch Chiều dài quỹ đạo là $2A = 8$ cm.
\itemch Tại $t=1$ s, $x = 4 \cos(2\pi - \pi/6) = 4 \cos(11\pi/6) = 4 . \sqrt{3}/2 = 2\sqrt{3}$ cm.
\end{itemchoice}
}
\end{ex}
\begin{ex}
Hình vẽ dưới đây mô tả dao động điều hòa của một vật.\\
\centerline{\begin{tikzpicture}[very thick,>=stealth']
\pgfplotsset{width=9cm,height=4cm,compat=1.4}
\begin{axis}[
axis lines=middle, xmin=0, xmax=2.2*pi, xstep=1,
ymin=-4.5, ymax=4.5, ystep=0.1,
grid=major,%Vẽ chế độ lưới theo trục tọa độ Oxy,
x grid style={dashed, black},
axis line style={very thick,Mapcolor}, % <--- dòng thêm vào
xlabel=$t~(ms)$,
xlabel style={above},
xtick ={2,4,6},
xticklabels={},
y grid style={dashed, black},
extra y ticks={0},
extra y tick label={0},
ylabel=$x~(cm)$,
ylabel style={above},
ytick={-4, 0,  4},
yticklabels={-4,0,4},]
\addplot[line width=1.2pt, domain=0:2*pi, samples=400, Mapcolor]{4*cos(deg(0.5*pi*x+pi/2))};
\addplot[Mapcolor,mark=*,mark options={fill=white,mark size=1.5pt}, nodes near coords={0,2}]
coordinates {(2, -4)};
\addplot[Mapcolor,mark=*,mark options={fill=white,mark size=1.5pt}, nodes near coords={0,4}]
coordinates {(4, -4)};
\addplot[Mapcolor,mark=*,mark options={fill=white,mark size=1.5pt}, nodes near coords={0,6}]
coordinates {(6, -4)};
\end{axis}
\end{tikzpicture}}
\choiceTF[t]
{\True Đồ thị dao động theo thời gian là một đường hình $\sin$}
{Biên độ dao động của vật là $2\ m$}
{Vật ở tại vị trí biên âm vào thời điểm $0,4 \ s$}
{\True Chu kì của dao động là $0,4\ s$}
\loigiai{
\begin{itemchoice}
\itemch Dao động điều hòa có đồ thị là hàm sin hoặc cosin.
\itemch Biên độ là giá trị cực đại của li độ, từ đồ thị thấy A = 4 cm.
\itemch Chu kì T = 0.8 s. Tại t = 0.4 s = T/2 vật ở biên âm x = -4 cm.
\itemch Tại t = 0, x = 0 và đang đi theo chiều âm, $\varphi = \pi/2$.
\end{itemchoice}
}
\end{ex}
\setcounter{ex}{0}
\PhanIII
\begin{ex}
Một chất điểm dao động điều hoà với phương trình $x=5 \cos \left(10 \pi t+\dfrac{\pi}{3}\right)\ cm$. Li độ của chất điểm khi pha dao động bằng $\pi$ là bao nhiêu centimet?
\shortans[oly]{-5}
\loigiai{
$x=5 \cos \pi=-5 \ cm$.
}
%\haidong{2}
\end{ex} \haidong{15}
\begin{ex}
Một chất điểm dao động điều hòa dọc theo trục $Ox$ với phương trình $x=2 \cos \left(2 \pi t-\dfrac{\pi}{6}\right)\ cm$. Li độ của chất điểm tại thời điểm $t=0,25 \ s$ bằng bao nhiêu centimet?
\shortans[oly]{1}
\loigiai{
$x=2 \cos \left(2 \pi . 0,25-\dfrac{\pi}{6}\right)=1 \ cm$
}
%\haidong{2}
\end{ex} \haidong{15}

\begin{ex}
Một chất điểm dao động điều hòa với phương trình $x=10 \cos \left(2 \pi t+\dfrac{\pi}{6}\right)\ cm$. Chu kì dao động của chất điểm bằng bao nhiêu giây?
\shortans[oly]{1}
\loigiai{
$T=\dfrac{2 \pi}{\omega}=\dfrac{2 \pi}{2 \pi}=1 \ s$
}
%\haidong{2}
\end{ex} \haidong{15}

\begin{ex}
Một chất điểm dao động có phương trình $x=10 \cos (15 t+\pi)$ (x tính bằng cm, t tính bằng ). Chất điểm này dao động với tần số góc bằng bao nhiêu rad/s?
\shortans[oly]{15}
\loigiai{
$\omega=15\ rad/s$
}
%\haidong{2}
\end{ex} \haidong{15}
\begin{ex}
Một chất điểm dao động điều hòa trên trục $Ox$ có phương trình $x=8 \cos \left(\pi t+\dfrac{\pi}{2}\right)\ cm$. Pha dao động của chất điểm khi $t=1 \ s$ bằng bao nhiêu rad (làm tròn kết quả đến chữ số hàng phần trăm)?
\shortans[oly]{4,71}
\loigiai{
4,71 rad.
}
%\haidong{2}
\end{ex} \haidong{15}
\begin{ex}
Một chất điểm dao động điều hoà có phương trình li độ theo thời gian là $x=5 \sqrt{3} \cos \left(10 \pi t+\dfrac{\pi}{3}\right)\ cm$. Tần số của dao động bằng bao nhiêu $Hz$?
\shortans[oly]{5}
\loigiai{
$f=\dfrac{\omega}{2 \pi}=\dfrac{10 \pi}{2 \pi}=5 \ Hz$.
}
%\haidong{1}
\end{ex} \haidong{15}
